\documentclass{article}
\usepackage[utf8]{inputenc}
\usepackage{geometry}
 \geometry{
 a4paper,
 total={170mm,257mm},
 left=20mm,
 top=20mm,
 }
 \usepackage{graphicx}
 \usepackage{titling}
 \usepackage{booktabs}
 \usepackage{listings}
 \usepackage{xcolor}
 \usepackage{hyperref}
 \usepackage{fancyvrb}
 \usepackage{float}
 \usepackage{amsmath}
 \definecolor{darkblue}{rgb}{0.0, 0.0, 0.5}

\RecustomVerbatimEnvironment{verbatim}{Verbatim}{formatcom=\color{darkblue}}

\setlength{\parindent}{0pt} %lui dau dong


\title{\textbf{DEEP LEARNING} \\ 
{\Large \textbf{Labwork 2 - Linear Regression}}}

\date{May 2026}
 
 \usepackage{fancyhdr}
\fancypagestyle{plain}{
    \fancyhf{} 
    \fancyfoot[L]{\thedate}
    \fancyfoot[R]{\thepage}
    \fancyhead[L]{University of Science and Technology of Hanoi}
    \fancyhead[R]{Deep Learning}
}
\makeatletter
\def\@maketitle{%
  \newpage
  % \null
  \vskip 1em %
  \begin{center}%
  \let \footnote \thanks
    {\LARGE \@title \par}%
    \vskip 1em%
    %{\large \@date}%
  \end{center}%
  \par
  \vskip 1em}
\makeatother

\usepackage{cmbright}

% Configuration for code blocks
\definecolor{codegreen}{rgb}{0,0.6,0}
\definecolor{backcolour}{rgb}{0.95,0.95,0.92}
\lstset{
    backgroundcolor=\color{backcolour},   
    commentstyle=\color{codegreen},
    basicstyle=\ttfamily\footnotesize,
    breaklines=true,
    showstringspaces=false
}

\begin{document}

\pagestyle{plain}
\maketitle

\noindent\begin{tabular}{@{}ll}
    Student: & Foo Cheung Nicolas - ES2540023\\
\end{tabular}

\section{Algorithm Implementation}

In this labwork, we implemented linear regression from scratch using gradient descent optimization. The goal is to learn parameters $w_0$ and $w_1$ such that the model:

\[
\hat{y} = w_1 x + w_0
\]

best fits the given dataset.


\subsection{Model Definition}

The linear regression model is defined as:

\[
f(x) = w_1 x + w_0
\]

\subsection{Loss Function}

We use the Mean Squared Error (MSE) with a factor $\frac{1}{2}$ for convenience in derivative computation:

\[
J(w_0, w_1) = \frac{1}{n} \sum_{i=1}^{n} \frac{1}{2} (\hat{y}_i - y_i)^2
\]

\subsection{Gradient Descent}

The parameters are updated using gradient descent:

\[
w_0 := w_0 - r \frac{\partial J}{\partial w_0}
\]
\[
w_1 := w_1 - r \frac{\partial J}{\partial w_1}
\]

where $r$ is the learning rate.

The partial derivatives are:

\[
\frac{\partial J}{\partial w_0} = \frac{1}{n} \sum (\hat{y}_i - y_i)
\]
\[
\frac{\partial J}{\partial w_1} = \frac{1}{n} \sum (\hat{y}_i - y_i)x_i
\]


\section{Effect of Learning Rate}

The learning rate $r$ significantly affects convergence:

\begin{itemize}
    \item \textbf{Small learning rate ($r = 0.00001$):} 
    The algorithm converges slowly but steadily toward the optimal solution.
    
    \item \textbf{Moderate learning rate ($r = 0.0005$):}
    Faster convergence with stable updates.

    \item \textbf{Large learning rate ($r = 0.001$):}
    The algorithm may diverge, causing the loss to increase rapidly and parameters to become unstable.
\end{itemize}

In this experiment, a moderate learning rate ($r = 0.0008$) was chosen to ensure stable convergence.

\section{Result Visualization}

The final model is plotted together with the dataset points. The learned regression line approximates the relationship between $x$ and $y$.

\begin{figure}[H]
    \centering
    \includegraphics[width=0.7\linewidth]{image.png}
    \caption{Linear regression plot}
    \label{fig:placeholder}
\end{figure}

\end{document}